\documentclass[11pt]{article}

\usepackage[letterpaper,margin=1in]{geometry}
\usepackage[T1]{fontenc}
\usepackage[utf8]{inputenc}
\usepackage{lmodern}
\usepackage{microtype}
\usepackage{amsmath,amssymb,mathtools,bm}
\usepackage{enumitem}
\usepackage{xcolor}
\usepackage{booktabs}
\usepackage{array}
\usepackage{hyperref}
\usepackage[most]{tcolorbox}
\usepackage{fancyhdr}
\usepackage{titlesec}
\usepackage{parskip}
\usepackage{needspace}

\definecolor{qftblue}{HTML}{245A8D}
\definecolor{qftgreen}{HTML}{2F6B4F}
\definecolor{qftred}{HTML}{A23B3B}
\definecolor{qftgray}{HTML}{F4F6F8}

\hypersetup{
  colorlinks=true,
  linkcolor=qftblue,
  urlcolor=qftblue,
  citecolor=qftgreen,
  pdftitle={PHYS 662 Lecture 9: Spin, Tensor Products, and Quantum Symmetries},
  pdfauthor={Nima Lashkari}
}

\pagestyle{fancy}
\fancyhf{}
\lhead{PHYS 662: Quantum Field Theory I}
\rhead{Lecture 9}
\cfoot{\thepage}
\setlength{\headheight}{14pt}

\titleformat{\section}{\Large\bfseries\color{qftblue}}{\thesection.}{0.6em}{}
\titleformat{\subsection}{\large\bfseries\color{qftgreen}}{\thesubsection}{0.6em}{}
\titleformat{\subsubsection}{\normalsize\bfseries\color{qftred}}{\thesubsubsection}{0.6em}{}

\setlist[itemize]{leftmargin=2em,itemsep=0.25em,topsep=0.4em}
\setlist[enumerate]{leftmargin=2em,itemsep=0.25em,topsep=0.4em}

\newcommand{\dd}{\mathrm{d}}
\newcommand{\R}{\mathbb{R}}
\newcommand{\Lag}{\mathcal{L}}
\newcommand{\cM}{\mathcal{M}}
\newcommand{\SO}{\mathrm{SO}}
\newcommand{\OO}{\mathrm{O}}

\newtcolorbox{keyidea}[1][]{
  enhanced, breakable, colback=qftgray, colframe=qftblue,
  boxrule=0.8pt, arc=1.5mm, title={#1}, fonttitle=\bfseries,
  coltitle=white
}
\newtcolorbox{remarkbox}[1][]{
  enhanced, breakable, colback=white, colframe=qftgreen,
  boxrule=0.7pt, arc=1.5mm, title={#1}, fonttitle=\bfseries,
  coltitle=white
}
\newtcolorbox{warningbox}[1][]{
  enhanced, breakable, colback=white, colframe=qftred,
  boxrule=0.7pt, arc=1.5mm, title={#1}, fonttitle=\bfseries,
  coltitle=white
}


\newcommand{\Th}{\Theta}
\newcommand{\Eul}{\mathcal E}
\newcommand{\Lie}{\mathscr L}

\newcommand{\C}{\mathbb C}
\newcommand{\Z}{\mathbb Z}
\newcommand{\SU}{\mathrm{SU}}
\newcommand{\UU}{\mathrm{U}}
\newcommand{\GL}{\mathrm{GL}}
\newcommand{\SL}{\mathrm{SL}}
\newcommand{\Sp}{\mathrm{Sp}}
\newcommand{\tr}{\operatorname{tr}}
\newcommand{\Ad}{\operatorname{Ad}}
\newcommand{\ad}{\operatorname{ad}}
\newcommand{\lieg}{\mathfrak g}
\newcommand{\lieh}{\mathfrak h}
\newcommand{\su}{\mathfrak{su}}
\newcommand{\so}{\mathfrak{so}}
\newcommand{\slie}{\mathfrak{sl}}


\newcommand{\cH}{\mathcal H}
\newcommand{\Spin}{\operatorname{Spin}}
\newcommand{\id}{\mathrm{id}}
\begin{document}
\pagenumbering{gobble}
\begin{titlepage}
\centering
\vspace*{1.2in}
{\scshape\Large PHYS 662 --- Quantum Field Theory I\par}
\vspace{0.35in}
{\Huge\bfseries\color{qftblue} Lecture 9\par}
\vspace{0.15in}
{\LARGE\bfseries Spin, Tensor Products,\par}
\vspace{0.06in}
{\LARGE\bfseries and Quantum Symmetries\par}
\vspace{0.16in}
{\large\color{qftgreen} Multiplets, Quantum Rays, and Covering Groups\par}
\vspace{0.6in}
{\large Nima Lashkari\par}
{\large Fall 2026\par}
\vfill
\begin{minipage}{0.86\textwidth}
\small
Starting from the Lie algebra of rotations, we construct the spin
multiplets and decompose tensor products into irreducible pieces. We then
ask how these representations act on quantum states. Rays and covering
groups explain why half-integer spin is compatible with spatial rotations,
and prepare the projective-representation analysis of Lecture 10.

\medskip
Course page: \url{https://nimalashkari.com/quantum-field-theory-2026-part-i/}
\end{minipage}
\vspace{0.5in}
\end{titlepage}
\pagenumbering{roman}
\begingroup\small
\tableofcontents
\endgroup
\newpage
\pagenumbering{arabic}

\section*{Roadmap}
\addcontentsline{toc}{section}{Roadmap}
Lecture 8 gave us Lie brackets, invariant subspaces, complete reducibility,
and Schur's lemma. We now use them to build representations and interpret
their physical meaning. By the end, we should be able to:
\begin{enumerate}
\item construct a spin multiplet from a highest-weight state;
\item decompose tensor products into multiplets with definite spin;
\item explain how quantum rays and the cover \(\SU(2)\to\SO(3)\)
allow half-integer spin.
\end{enumerate}
We first solve the finite-dimensional algebra problem, then ask which
global group those matrices represent. Lecture 10 will study the general
phase-multiplication law and the algebra of conserved charges.

\begin{keyidea}[Tools carried forward from Lecture 8]
We use \(U=e^{-i\theta^aT_a}\), \(\hbar=1\), and
\([J_i,J_j]=i\epsilon_{ijk}J_k\), with \(\epsilon_{123}=+1\).
Representation spaces are complex unless stated otherwise. Compactness
allows a finite-dimensional representation to be made unitary, so its
generators are Hermitian. In an irreducible complex representation,
Schur's lemma makes every operator commuting with the whole action a
scalar multiple of the identity.
\end{keyidea}

\section{Rotations and the spin multiplets}
\label{sec:spin}
\subsection{The two-dimensional representation of \texorpdfstring{\(\SU(2)\)}{SU(2)}}
The Pauli matrices are
\[
 \sigma_1=\begin{pmatrix}0&1\\1&0\end{pmatrix},\qquad
 \sigma_2=\begin{pmatrix}0&-i\\i&0\end{pmatrix},\qquad
 \sigma_3=\begin{pmatrix}1&0\\0&-1\end{pmatrix}.
\]
Their multiplication law,
\(\sigma_i\sigma_j=\delta_{ij}I+i\epsilon_{ijk}\sigma_k\), gives
\begin{equation}
 J_i=\frac{\sigma_i}{2},\qquad
 [J_i,J_j]=i\epsilon_{ijk}J_k,\qquad
 \tr(J_iJ_j)=\frac12\delta_{ij}.
 \label{eq:pauli-generators}
\end{equation}
The factors of two are essential. The Pauli matrices themselves have
commutator \([\sigma_i,\sigma_j]=2i\epsilon_{ijk}\sigma_k\).
For a unit vector \(\bm n\), \((\bm n\cdot\bm\sigma)^2=I\), so
\begin{equation}
 U(\theta,\bm n)
 =e^{-i\theta\bm n\cdot\bm\sigma/2}
 =\cos\frac\theta2\,I-i\sin\frac\theta2\,\bm n\cdot\bm\sigma.
 \label{eq:su2-exponential}
\end{equation}
These matrices are unitary with determinant one. They display the
three real parameters of the group, while acting on a two-dimensional
complex vector space. The group dimension and representation dimension
are different quantities.

\Needspace{9\baselineskip}
\subsection{Deriving the spin representations}
The aim is to find every possible finite spin multiplet without guessing
matrices of each size. We diagonalize one generator and use the other two
to move between its eigenvectors. Compactness lets us choose a unitary
representation, so the generators are Hermitian. In an irreducible
representation, define
\[
 J_\pm=J_1\pm iJ_2,\qquad
 J^2=J_1^2+J_2^2+J_3^2.
\]
The commutation relations imply
\begin{equation}
 [J_3,J_\pm]=\pm J_\pm,\qquad
 [J_+,J_-]=2J_3,\qquad [J^2,J_i]=0.
 \label{eq:ladder-algebra}
\end{equation}
By Schur's lemma \(J^2=cI\). Since \(J_3\) is Hermitian, choose a
normalized eigenvector with eigenvalue \(m\). The raising relation means
\[
 J_3(J_+|m\rangle)
 =(J_+J_3+J_+)|m\rangle=(m+1)J_+|m\rangle.
\]
Thus \(J_+|m\rangle\) either vanishes or has eigenvalue \(m+1\);
\(J_-\) similarly lowers m. A \emph{weight} is this eigenvalue, and a
\emph{highest-weight vector} is the state at the top of the ladder.
Since \(J_+^\dagger=J_-\), its norm is computed as
\(\|J_+|m\rangle\|^2=\langle m|J_-J_+|m\rangle\). The identities
\[
 J_-J_+=J^2-J_3(J_3+1),\qquad
 J_+J_-=J^2-J_3(J_3-1)
\]
therefore give the squared norms
\begin{equation}
 \|J_+|m\rangle\|^2=c-m(m+1),\qquad
 \|J_-|m\rangle\|^2=c-m(m-1).
 \label{eq:ladder-norms}
\end{equation}
Let the largest eigenvalue be \(j\). The raising operator must annihilate
its eigenvector, so \(c=j(j+1)\). If \(m_{\min}\) is the smallest
eigenvalue, the second norm gives
\(m_{\min}(m_{\min}-1)=j(j+1)\). Its roots are \(-j\) and \(j+1\);
the latter cannot be a lowest eigenvalue below \(j\). Hence
\(m_{\min}=-j\). The number of lowering steps is the nonnegative
integer \(2j\).

Choosing phases consistently, we obtain
\begin{align}
 j&=0,\tfrac12,1,\tfrac32,\ldots,
 & \dim V_j&=2j+1,\label{eq:spin-labels}\\
 J_3|j,m\rangle&=m|j,m\rangle,
 &m&=-j,-j+1,\ldots,j,\\
 J^2|j,m\rangle&=j(j+1)|j,m\rangle,\\
 J_\pm|j,m\rangle
 &=\sqrt{(j\mp m)(j\pm m+1)}\,|j,m\pm1\rangle.
 \label{eq:ladder-action}
\end{align}
The chain generated from the highest-weight vector is invariant under
all three generators. Irreducibility makes it the entire representation;
conversely, the displayed matrices construct an irreducible representation
for each allowed \(j\). Thus this is the complete finite-dimensional
complex irreducible list for \(\SU(2)\).

For \(j=0\), the generators are zero and every group element acts as the
one-by-one matrix \(1\). The trivial group representation is the identity,
whereas its infinitesimal representation is zero. For \(j=\tfrac12\),
Eq.~\eqref{eq:ladder-action} reproduces \(J_i=\sigma_i/2\).

\Needspace{16\baselineskip}
\subsection{Spin one in two bases}
To turn the ladder formula into matrices, remember that the kth column
contains the image of the kth basis vector. For j = 1,
\(J_+|1,0\rangle=\sqrt2|1,1\rangle\) and
\(J_+|1,-1\rangle=\sqrt2|1,0\rangle\), while
\(J_+|1,1\rangle=0\). Then use
\(J_1=(J_++J_-)/2\) and \(J_2=(J_+-J_-)/(2i)\).
In the ordered basis \((|1,1\rangle,|1,0\rangle,|1,-1\rangle)\),
\begin{equation}
 \begin{split}
 J_1&=\frac1{\sqrt2}
 \begin{pmatrix}0&1&0\\1&0&1\\0&1&0\end{pmatrix},\qquad
 J_2=\frac1{\sqrt2}
 \begin{pmatrix}0&-i&0\\i&0&-i\\0&i&0\end{pmatrix},\\
 J_3&=\begin{pmatrix}1&0&0\\0&0&0\\0&0&-1\end{pmatrix}.
 \end{split}
 \label{eq:spin-one-matrices}
\end{equation}
They obey \([J_i,J_j]=i\epsilon_{ijk}J_k\) and \(J^2=2I_3\).
The operators \(J_\pm\), in contrast, are triangular and are not
Hermitian generators of real one-parameter rotation subgroups.

In a Cartesian basis of \(\C^3\), the same representation has generators
\begin{equation}
 (L_i)_{jk}=-i\epsilon_{ijk}.
 \label{eq:cartesian-spin-one}
\end{equation}
The matrices \(e^{-i\theta L_i}\) are real spatial rotations. For example,
\[
 e^{-i\theta L_3}=
 \begin{pmatrix}\cos\theta&-\sin\theta&0\\
 \sin\theta&\cos\theta&0\\0&0&1\end{pmatrix}.
\]
An explicit change from Cartesian to the magnetic basis is supplied by
\[
 |1,1\rangle=-\frac{\bm e_x+i\bm e_y}{\sqrt2},\qquad
 |1,0\rangle=\bm e_z,\qquad
 |1,-1\rangle=\frac{\bm e_x-i\bm e_y}{\sqrt2}.
\]
If \(V\) has these three vectors as columns, then \(V\) is unitary and
\(V^\dagger L_iV=J_i\). Different-looking generator matrices can therefore
be equivalent descriptions of the same representation.

\section{Tensor products and irreducible tensors}
\label{sec:tensors}
\subsection{Direct sums versus tensor products}
Two constructions recur in field theory. The direct sum \(V\oplus W\)
puts two multiplets side by side:
\[
 (R\oplus S)(g)=
 \begin{pmatrix}R(g)&0\\0&S(g)\end{pmatrix},\qquad
 \dim(V\oplus W)=\dim V+\dim W.
\]
The tensor product describes pairs of basis labels. Two two-component
spaces, for example, have the four product basis vectors
\(e_1\otimes f_1,e_1\otimes f_2,e_2\otimes f_1,e_2\otimes f_2\).
A general tensor is a linear combination of these products; it need not
itself be a single product. The group acts on both labels:
\begin{equation}
 (R\otimes S)(g)=R(g)\otimes S(g),\qquad
 \dim(V\otimes W)=(\dim V)(\dim W).
 \label{eq:tensor-product}
\end{equation}
If \(v\mapsto R(g)v\) and \(w\mapsto S(g)w\), then
\(v\otimes w\mapsto R(g)v\otimes S(g)w\).
Differentiating the group action gives
\begin{equation}
 T_a^{V\otimes W}=T_a^V\otimes I_W+I_V\otimes T_a^W.
 \label{eq:tensor-generators}
\end{equation}
This is the familiar addition rule for angular momentum.

The tensor product is often reducible. Decomposing it means finding a
\emph{direct sum} of invariant subspaces, so the symbol changes from
\(\otimes\) to \(\oplus\). The group and the field of scalars should also
be specified. For example, the real two-dimensional defining
representation of \(\SO(2)\) is irreducible over \(\R\); over \(\C\),
it splits into the one-dimensional weights \(+1\) and \(-1\).
In the spin decompositions below we use complex vector spaces.

\subsection{Conjugate and dual representations}
If a column \(v\) transforms by R(g), then its complex conjugate
transforms by \(R(g)^*\). A dual covector \(\ell\) instead transforms as
\(\ell\mapsto\ell R(g)^{-1}\), so the pairing \(\ell(v)\) is invariant.
On columns of dual components this is the representation
\(R^\vee(g)=R(g)^{-T}\). For a unitary representation,
\(R(g)^{-T}=R(g)^*\); an invariant Hermitian inner product identifies
the conjugate and dual representations.

For a fundamental U(N) multiplet this becomes
\begin{equation}
 \Phi'^A=U^A{}_B\Phi^B,\qquad
 \Phi_A'^\dagger=\Phi_B^\dagger(U^\dagger)^B{}_A,
 \qquad \Phi'^\dagger\Phi'=\Phi^\dagger\Phi.
 \label{eq:fundamental-dual}
\end{equation}
An upper or lower internal index here records the representation in which
it transforms; it is not a spacetime index raised or lowered with the
Minkowski metric. The pairing is invariant, but the field norm is not
constrained to equal one. Lecture 10 will use this transformation law to
construct conserved currents.

\subsection{A rank-two tensor under spatial rotations}
Let \(V_1\) denote the three-dimensional spin-one representation.
A general tensor \(M_{ij}\) belongs to \(V_1\otimes V_1\), a
nine-dimensional space, and transforms as
\begin{equation}
 M'_{ij}=R_{ia}R_{jb}M_{ab},\qquad M'=RMR^T,
 \qquad R\in\SO(3).
 \label{eq:rank-two-action}
\end{equation}
Every such tensor has a unique decomposition
\begin{equation}
 \boxed{M_{ij}=\frac{\tr M}{3}\delta_{ij}+A_{ij}+Q_{ij}},
 \label{eq:tensor-split}
\end{equation}
where
\begin{align}
 A_{ij}&=\tfrac12(M_{ij}-M_{ji}),\label{eq:antisymmetric-part}\\
 Q_{ij}&=\tfrac12(M_{ij}+M_{ji})-\tfrac13\delta_{ij}\tr M,
 \qquad Q_{ij}=Q_{ji},\quad Q_{ii}=0.
 \label{eq:traceless-part}
\end{align}
These are invariant subspaces: transposition commutes with
\(M\mapsto RMR^T\), \(R\delta R^T=\delta\), and
\(\tr(RMR^T)=\tr M\). Their dimensions are
\[
 \underbrace{1}_{\text{trace}}+
 \underbrace{3}_{\text{antisymmetric}}+
 \underbrace{5}_{\text{symmetric traceless}}=9.
\]
The trace is a rotational scalar, hence spin zero. The antisymmetric
piece is equivalent to a vector through
\begin{equation}
 a_k=\tfrac12\epsilon_{kij}A_{ij},\qquad
 A_{ij}=\epsilon_{ijk}a_k.
 \label{eq:axial-dual}
\end{equation}
The invariance of \(\epsilon_{ijk}\) under proper rotations gives
\(a'_k=R_{k\ell}a_\ell\), so this piece is spin one. Under an improper
orthogonal transformation it instead acquires an extra factor
\(\det R\); it is then called an axial vector.

To identify the symmetric traceless part, take the tensor product of
two spin-one highest-weight vectors. It is symmetric, has weight two,
and has zero trace, because the Cartesian vector
\(\bm e_x+i\bm e_y\) has zero contraction with itself without complex
conjugation. Repeated lowering generates a five-dimensional spin-two
subspace. Since the entire symmetric traceless space has dimension five,
it is precisely this irreducible representation. We conclude
\begin{equation}
 \boxed{V_1\otimes V_1\cong V_0\oplus V_1\oplus V_2},
 \qquad \bm3\otimes\bm3=\bm1\oplus\bm3\oplus\bm5.
 \label{eq:three-times-three}
\end{equation}
Bold numbers in this shorthand label representations by their dimensions.
The corresponding numerical identity is \(3\times3=1+3+5\).

\begin{remarkbox}[A tensor split you can check entry by entry]
Take a matrix with only two nonzero entries:
\[
 M=\begin{pmatrix}1&2&0\\0&0&0\\0&0&0\end{pmatrix}
 =\frac13 I+
 \begin{pmatrix}0&1&0\\-1&0&0\\0&0&0\end{pmatrix}
 +\frac13\begin{pmatrix}2&3&0\\3&-1&0\\0&0&-1\end{pmatrix}.
\]
The pieces are respectively scalar, antisymmetric, and symmetric
traceless. The antisymmetric piece corresponds to \(\bm a=\bm e_z\).
A rotation mixes entries inside each of these three spaces, while
preserving their defining symmetry and trace properties. That is the
meaning of the irreducible decomposition.
\end{remarkbox}

\subsection{Higher rank and angular-momentum addition}
A symmetric rank-\(\ell\) tensor in three dimensions has
\(\binom{\ell+2}{2}\) independent components. For \(\ell\geq2\),
the trace map onto symmetric rank-\(\ell-2\) tensors is surjective.
One way to see this is to identify symmetric tensors with homogeneous
polynomials: trace is proportional to the Laplacian, which maps degree
\(\ell\) polynomials onto degree \(\ell-2\) polynomials. For a polynomial
\(q\) of degree \(k\),
\[
 \nabla^2(r^2q)=(4k+6)q+r^2\nabla^2q;
\]
iterating this identity with the successive Laplacians of \(q\)
constructs a preimage. Therefore the symmetric traceless subspace has
\begin{equation}
 \binom{\ell+2}{2}-\binom{\ell}{2}=2\ell+1
 \label{eq:higher-tensor-dimension}
\end{equation}
components. Its highest-weight tensor is
\((\bm e_x+i\bm e_y)^{\otimes\ell}\); it generates the spin-\(\ell\)
representation. Rank zero and rank one give the same count directly.
This is why spherical harmonics of degree \(\ell\) have
\(m=-\ell,\ldots,\ell\).

More generally, angular-momentum addition gives the Clebsch--Gordan rule
\begin{equation}
 V_{j_1}\otimes V_{j_2}
 \cong\bigoplus_{j=|j_1-j_2|}^{j_1+j_2}V_j,
 \label{eq:cg-rule}
\end{equation}
where the sum increases in steps of one. The highest weight of the
product is \(j_1+j_2\). Remove that irreducible multiplet, then repeat
with the largest remaining weight. This gives one copy of each spin in
the displayed range, and their dimensions sum to
\((2j_1+1)(2j_2+1)\). For two spin-one-half representations,
\[
 \bm2\otimes\bm2=\bm1\oplus\bm3.
\]
The antisymmetric combination
\((|{+}{-}\rangle-|{-}{+}\rangle)/\sqrt2\) is the singlet; the three
symmetric combinations form the triplet. For two spin-one representations,
Eq.~\eqref{eq:cg-rule} agrees with the explicit tensor decomposition above.

A rank-three Cartesian tensor has 27 components and transforms with
three factors of \(R\). Applying the same rule twice gives
\[
 V_1^{\otimes3}\cong V_0\oplus3V_1\oplus2V_2\oplus V_3,
 \qquad 27=1+3(3)+2(5)+7.
\]
Here \(3V_1\) means three copies in a direct sum. Only the fully symmetric
traceless rank-three part is the seven-dimensional spin-three multiplet;
a general rank-three tensor contains the other spins as well.

\subsection{Matrix fields and the adjoint of \texorpdfstring{\(\UU(N)\)}{U(N)}}
A matrix with one fundamental and one dual index transforms by
\begin{equation}
 M'^A{}_B=U^A{}_C M^C{}_D(U^\dagger)^D{}_B,
 \qquad M'=UMU^\dagger.
 \label{eq:matrix-field}
\end{equation}
Its trace is invariant and its traceless part stays traceless. Thus, for
\(N\geq2\),
\begin{equation}
 \bm N\otimes\overline{\bm N}
 \cong\bm1\oplus\operatorname{Adj}_{\SU(N)},\qquad
 N^2=1+(N^2-1),\qquad
 M=\frac{\tr M}{N}I+\left(M-\frac{\tr M}{N}I\right).
 \label{eq:fundamental-antifundamental}
\end{equation}
This is the complexified adjoint action of U(N). Its full dimension is
\(N^2=\dim\UU(N)\), but it is reducible: the identity matrix is a singlet
and the traceless subspace is the SU(N) adjoint. The central U(1) subgroup
acts trivially by conjugation, so the U(N) adjoint is not faithful either.

To obtain the generator matrices, choose a Hermitian basis \(t_a\) and define
\(Ut_bU^\dagger=t_aD(U)^a{}_b\). Expanding
\(U=e^{-i\epsilon^ct_c}\) gives
\begin{equation}
 D(U)^a{}_b=\delta^a{}_b+\epsilon^cf_{cb}{}^a+O(\epsilon^2),
 \qquad (t_c^{\rm ad})^a{}_b=i f_{cb}{}^a,
 \qquad D(U)=e^{-i\epsilon^ct_c^{\rm ad}}.
 \label{eq:adjoint-matrices}
\end{equation}
For a basis with \(\tr(t_at_b)=\delta_{ab}/2\), the fully lowered
structure constants are antisymmetric, and the same formula reads
\((t_c^{\rm ad})_{ab}=-if_{cab}\). The U(1) generator has zero adjoint
matrix. This separates the dimension of the adjoint representation from
whether it is irreducible or faithful, and prepares the transformation
law for the currents in Lecture 10.

\section{Quantum symmetries act on rays}
\label{sec:rays}
\subsection{Wigner's theorem}
A pure state is a ray in a complex Hilbert space \(\cH\): two nonzero
vectors related by a nonzero complex factor describe the same state.
For normalized vectors the redundancy is a phase, and the transition
probability is
\[
 P(\psi,\chi)=|\langle\psi|\chi\rangle|^2.
\]
For example, let \(\lvert0\rangle,\lvert1\rangle\) be orthonormal and set
\(\lvert\pm\rangle=(\lvert0\rangle\pm\lvert1\rangle)/\sqrt2\).
Multiplying \(\lvert+\rangle\) by one common phase leaves its ray
unchanged, whereas changing the relative sign gives the orthogonal
state \(\lvert-\rangle\). Only the common phase is redundant.

Wigner's theorem states that a bijection of rays preserving these
probabilities is implemented by a unitary or antiunitary operator,
unique up to a phase. The physical input is preservation of probabilities,
not a demand that the complex overlap itself remain unchanged.

A unitary operator is complex linear and satisfies
\(\langle U\psi|U\chi\rangle=\langle\psi|\chi\rangle\).
An antiunitary operator \(A\) is conjugate linear:
\begin{equation}
 A(c_1\psi+c_2\chi)=c_1^*A\psi+c_2^*A\chi,
 \qquad \langle A\psi|A\chi\rangle=\langle\chi|\psi\rangle.
 \label{eq:antiunitary}
\end{equation}
Both preserve transition probabilities. In a chosen orthonormal basis,
\(A=VK\), where \(V\) is unitary and \(K\) conjugates coefficients.
Thus an antiunitary map is not just a complex matrix acting linearly.
Continuous transformations connected to the identity have unitary
implementers; disconnected transformations can be antiunitary.

For a spin-one-half degree of freedom, time reversal can be implemented
by \(\mathcal T=MK\), with \(M=i\sigma_2\). Since
\(KM=M^*K\) and \(K^2=I\), its square is \(MM^*=-I\).
Using the Pauli matrices also gives
\[
 \mathcal T\bm J\mathcal T^{-1}=-\bm J,\qquad
 \mathcal T^2=-I,\qquad \bm J=\bm\sigma/2.
\]
The minus sign in its square acts trivially on a ray. This example also
shows why replacing an antiunitary operator by its matrix factor alone
would give the wrong transformation law.

\subsection{Field components and quantum states}
A Lorentz vector field has four components transformed by a
finite-dimensional matrix preserving the Minkowski metric, generally not
a positive-definite inner product. The corresponding quantum symmetry
acts unitarily on a Hilbert space that includes continuous momentum
labels. These are different representations on different spaces.

Even a scalar's trivial component representation does not make its
function-space representation trivial:
\[
 \phi'(x)=\phi\bigl(\Lambda^{-1}(x-a)\bigr).
\]
Translations act by shifts of the argument; their infinitesimal generators
are derivatives. Infinite-dimensional group operators are not in general
differential operators themselves.

A stable relativistic particle's one-particle space carries a
positive-energy irreducible unitary representation of the appropriate
Poincar\'e cover. This concerns all its momentum states and applies also
to stable composite particles; it does not define microscopic
fundamentality. The present lecture supplies the group-theoretic reasons
that a cover may be needed.

\subsection{Why multiplication can hold only up to phase}
Choose one unitary operator for each symmetry connected to the identity.
Two successive symmetries and their group product have the same action
on rays, so their operators can differ by a phase:
\begin{equation}
 U(g)U(h)=e^{i\omega(g,h)}U(gh).
 \label{eq:projective-preview}
\end{equation}
This is called a projective representation. The phase is one scalar on
the entire representation space. In this lecture we need this idea to
interpret the sign of a \(2\pi\) rotation. Lecture 10 will derive the
consistency condition on \(\omega\) and study which phases can be removed.

\section{Covering groups and the global origin of spin}
\label{sec:covers}
\subsection{Loops, the universal cover, and the circle}
A connected group is path connected as a manifold. It is simply connected
if every closed loop contracts to a point. The fundamental group
\(\pi_1(G)\) records loop classes, with multiplication given by
concatenation. The universal covering group \(p:\widetilde G\to G\)
is connected and simply connected; locally \(p\) is an isomorphism.
For a connected Lie group its kernel is discrete and central and is
isomorphic to \(\pi_1(G)\).

For example,
\begin{equation}
 \widetilde{\SO(2)}=\R,
 \qquad p(\theta)=R(\theta),\qquad
 \ker p=2\pi\Z,\qquad \pi_1(\SO(2))\cong\Z.
 \label{eq:circle-cover}
\end{equation}
The covering coordinate remembers how many complete turns have occurred.
The circle and the line have the same one-dimensional Lie algebra, but
different allowed ordinary characters:
\[
 R_q(t)=e^{-iqt}\ (q\in\R)\quad\text{on }\R;
 \qquad R_q(\theta)=e^{-iq\theta}\ (q\in\Z)\quad\text{on }\UU(1).
\]
The integer condition is exactly \(R_q(\theta+2\pi)=R_q(\theta)\).
For the finite cyclic group \(\Z_N\), the characters instead are
\[
 \chi_q(k)=e^{2\pi i qk/N},\qquad k,q=0,\ldots,N-1.
\]
The characters are fixed by the image of a single generator: its Nth
power must be one. Each finite group has zero-dimensional Lie algebra.

A finite-dimensional Lie algebra representation integrates uniquely to
its connected simply connected group. To \emph{descend} means to define
the action using g alone, without choosing a lift \(\widetilde g\).
Two lifts differ by \(k\in\ker p\), and
\[
 \widetilde U(\widetilde g k)
 =\widetilde U(\widetilde g)\widetilde U(k).
\]
They give the same operator if \(\widetilde U(k)=I\), and the same
action on rays if \(\widetilde U(k)\) is a scalar phase. These are the
conditions for ordinary and projective descent, respectively.
For an irreducible unitary
representation, central elements act as scalars by Schur's lemma;
a reducible representation requires checking the scalar condition on the
whole space.

\subsection{The double cover of spatial rotations}
The matrices in Eq.~\eqref{eq:su2-exponential} already display a
half-angle. We now explain how their two-dimensional action is related
to a three-dimensional spatial rotation.
Associate \(\bm v\in\R^3\) with the traceless Hermitian matrix
\[
 V=\bm v\cdot\bm\sigma
 =\begin{pmatrix}v_z&v_x-iv_y\\v_x+iv_y&-v_z\end{pmatrix}.
\]
Conjugation defines a real linear map \(R_U\) by
\begin{equation}
 UVU^\dagger=(R_U\bm v)\cdot\bm\sigma.
 \label{eq:su2-cover}
\end{equation}
It preserves \(|\bm v|^2=\tfrac12\tr V^2\), so it is orthogonal, and
has determinant one by continuity. For a rotation about the z-axis,
\[
 U_z=\begin{pmatrix}e^{-i\theta/2}&0\\0&e^{i\theta/2}\end{pmatrix},
 \qquad
 \begin{pmatrix}v'_x\\v'_y\\v'_z\end{pmatrix}
 =\begin{pmatrix}\cos\theta&-\sin\theta&0\\
 \sin\theta&\cos\theta&0\\0&0&1\end{pmatrix}
 \begin{pmatrix}v_x\\v_y\\v_z\end{pmatrix}.
\]
The general axis follows from the Pauli multiplication law. Thus every
proper rotation is obtained. The kernel consists of matrices commuting
with all Pauli matrices, namely \(\{I,-I\}\), and
\begin{equation}
 \boxed{\SO(3)\cong\SU(2)/\{I,-I\}},\qquad
 \so(3)\cong\su(2).
 \label{eq:double-cover}
\end{equation}
Every SU(2) matrix can be written \(a_0I-i\bm a\cdot\bm\sigma\), with
real coefficients satisfying \(a_0^2+|\bm a|^2=1\). Hence SU(2) is the
three-sphere, which is simply connected. It is the universal cover, and
\begin{equation}
 \pi_1(\SU(2))=0,\qquad \pi_1(\SO(3))=\Z_2.
 \label{eq:fundamental-groups}
\end{equation}
A loop describing a 2\(\pi\) rotation lifts to a path from I to -I;
traversing it twice gives a closed loop on the cover, which can contract.

\subsection{Which spin representations descend?}
The spin-j multiplet has \(m=-j,-j+1,\ldots,j\), so
\begin{equation}
 e^{-2\pi iJ_3}=(-1)^{2j}I.
 \label{eq:two-pi}
\end{equation}
For integer j, the kernel acts trivially and the representation descends
to an ordinary representation of SO(3). For half-integer j, it acts by
minus the identity: the representation is ordinary on SU(2) and
projective on SO(3). Locally the same commutator algebra holds in both
cases: \([J_i,J_j]=i\epsilon_{ijk}J_k\). The half-integer possibility
comes from the global group, with no extra identity term in this bracket.

The spin-one representation of SU(2) is irreducible but not faithful:
-I acts trivially. The spin-one-half representation is both irreducible
and faithful. These facts concern a particular action on a vector space,
not whether the group has normal subgroups.

A direct sum can fail even projective descent. On
\(V_0\oplus V_{1/2}\), the kernel element -I acts as
\(\operatorname{diag}(1,-I_2)\), which is not one scalar on the entire
space. Consequently not every ordinary representation of the covering
group defines a projective representation of the base group on all its
rays. The common scalar action of the kernel is the essential condition.

\subsection{Higher-dimensional spin groups}
For \(n\geq3\), the universal cover of SO(n) is its double cover Spin(n).
Some useful low-dimensional identifications are
\begin{equation}
 \begin{split}
 \Spin(3)&\cong\SU(2),\qquad
 \Spin(4)\cong\SU(2)\times\SU(2),\\
 \Spin(5)&\cong\Sp(2)\cong\mathrm{USp}(4),\qquad
 \Spin(6)\cong\SU(4).
 \end{split}
 \label{eq:spin-isomorphisms}
\end{equation}
Here Sp(2) is the compact symplectic group. Spin(2) is an exception to the
universal-cover statement: it is a circle that double covers SO(2),
whereas the universal cover of SO(2) is the line.

In Spin(4), the covering kernel is \(\{(I,I),(-I,-I)\}\). An irreducible
representation labeled \((j_L,j_R)\) descends to SO(4) precisely when
\(j_L+j_R\) is an integer. This illustrates how the global quotient adds
a condition beyond the direct-sum Lie algebra found in Lecture 8.

\clearpage
\section*{Questions and checks}
\addcontentsline{toc}{section}{Questions and checks}
\begin{enumerate}
\item Use Schur's lemma to explain why \(J^2\) has a single eigenvalue
in an irreducible complex spin representation. Compute it for spin two.
\item Verify the spin-one matrices in Eq.~\eqref{eq:spin-one-matrices},
including \(J^2=2I\) and the unitary change to the Cartesian basis.
\item Construct \(J_3\) and \(J_+\) for spin \(\tfrac32\) from the ladder
formula. Check the dimension and both endpoint conditions.
\item Verify that the three operations extracting the trace,
antisymmetric part, and symmetric traceless part in
Eq.~\eqref{eq:tensor-split} are mutually annihilating projectors whose
sum is the identity on the tensor space.
\item Compare the decompositions of \(M_{ij}=v_iw_j\) under SO(3) and
\(M^A{}_B=v^Aw_B^*\) under U(N). Identify their invariant scalar pieces
and explain why the U(N) adjoint is reducible.
\item Obtain the spin content of a general rank-three Cartesian tensor
using Eq.~\eqref{eq:cg-rule}, and check the total dimension. Compare it
with the fully symmetric rank-three tensor.
\item Verify that \(i\sigma_2K\) is antiunitary, reverses all three spin
operators, and squares to \(-I\). Why does its square fix every ray?
\item Compare the universal cover of SO(2), Spin(2), and SO(2) itself.
Derive the integer-weight condition for an ordinary circle character.
\item Compute \(U_zVU_z^\dagger\) explicitly and recover the ordinary
rotation of \((v_x,v_y)\). Identify the kernel of the full covering map.
\item Which of \(V_0\oplus V_1\), \(V_{1/2}\oplus V_{3/2}\), and
\(V_0\oplus V_{1/2}\) descend to ordinary or projective representations
of SO(3) on the whole vector space?
\end{enumerate}

\clearpage
\section*{Summary}
\addcontentsline{toc}{section}{Summary}
\begin{keyidea}[From an algebraic multiplet to a physical rotation]
The ladder calculation determines the allowed finite-dimensional
\(\SU(2)\) multiplets. Tensor products organize composite objects.
Quantum states are rays, so a scalar phase can be invisible on states.
The global cover then explains why both integer and half-integer spin
can describe rotations, even though only integer spin gives an ordinary
representation of \(\SO(3)\).
\end{keyidea}
\begin{itemize}
\item The spin-\(j\) multiplet has dimension \(2j+1\) and
\(J^2=j(j+1)I\). Ladder operators construct the matrices explicitly.
\item Tensor-product generators add. A rank-two spatial tensor splits
as \(\bm3\otimes\bm3=\bm1\oplus\bm3\oplus\bm5\): trace,
antisymmetric, and symmetric traceless pieces.
\item The dual action preserves a pairing. For U(N), a matrix field
transforms by conjugation and splits into a singlet and the SU(N)
adjoint; the central U(1) acts trivially on that matrix field.
\item Wigner's theorem implements probability-preserving ray symmetries
by unitary or antiunitary operators. Field-component and Hilbert-space
representations act on different spaces.
\item \(\SU(2)\) is the simply connected double cover of \(\SO(3)\).
The kernel element \(-I\) acts as \((-1)^{2j}\). Ordinary descent
requires the kernel to act as the identity; projective descent requires
one scalar on the whole space.
\item Lecture 10 studies the general projective phase, derives its
possible central Lie algebra term, and realizes symmetry generators as
Noether charges.
\end{itemize}

\section*{Further reading}
\addcontentsline{toc}{section}{Further reading}
These notes reorganize the supplied symmetry and representation slide
sets into Lectures 8--10. For mathematical background and further proofs:
\begin{itemize}
\item B. C. Hall, \emph{An Elementary Introduction to Groups and
Representations}, \href{https://arxiv.org/abs/math-ph/0005032}{arXiv:math-ph/0005032}.
\end{itemize}
\end{document}
